% docs/slides/preamble.tex -- 五份 deck 共享
\documentclass[aspectratio=169,11pt]{beamer}

% ==== 中文与字体 ====
\usepackage[UTF8,fontset=none]{ctex}
\usepackage{fontspec}

\IfFontExistsTF{PingFang SC}{%
  \setCJKmainfont{PingFang SC}[BoldFont=PingFang SC,ItalicFont=PingFang SC,AutoFakeSlant=0.15]
  \setCJKsansfont{PingFang SC}[ItalicFont=PingFang SC,AutoFakeSlant=0.15]
  \setCJKmonofont{PingFang SC}
}{%
  \IfFontExistsTF{Source Han Sans SC}{%
    \setCJKmainfont{Source Han Sans SC}
    \setCJKsansfont{Source Han Sans SC}
  }{%
    \IfFontExistsTF{Noto Sans CJK SC}{%
      \setCJKmainfont{Noto Sans CJK SC}
      \setCJKsansfont{Noto Sans CJK SC}
    }{%
      \setCJKmainfont{Heiti SC}
      \setCJKsansfont{Heiti SC}
    }
  }
}

\IfFontExistsTF{Helvetica Neue}{\setmainfont{Helvetica Neue}}{%
  \IfFontExistsTF{Inter}{\setmainfont{Inter}}{}}
\IfFontExistsTF{JetBrains Mono}{\setmonofont{JetBrains Mono}[Scale=0.85]}%
  {\setmonofont{Menlo}[Scale=0.85]}

% ==== 颜色 ====
\usepackage{xcolor}
\definecolor{cMain}{HTML}{1F4E79}
\definecolor{cAccent}{HTML}{E07A1F}
\definecolor{cGray}{HTML}{595959}
\definecolor{cMute}{HTML}{F4F4F4}
\definecolor{cOK}{HTML}{2E7D32}
\definecolor{cBad}{HTML}{B23A3A}
\definecolor{cWarn}{HTML}{C8A300}

% ==== 主题 ====
\usepackage{tcolorbox}
\usetheme{metropolis}
\metroset{progressbar=frametitle,numbering=fraction,block=fill}
\setbeamercolor{frametitle}{bg=cMain,fg=white}
\setbeamercolor{progress bar}{fg=cAccent,bg=cMute}
\setbeamercolor{alerted text}{fg=cAccent}
\setbeamercolor{title}{fg=cMain}

% metropolis 默认尝试 Fira Sans，未安装时 fallback 到 Latin Modern Sans。
% 这里在主题加载后显式覆盖，强制使用 macOS 上一定存在的 Helvetica Neue，
% 跟中文 PingFang 风格一致；找不到则按顺序尝试 Inter / Arial。
\IfFontExistsTF{Helvetica Neue}{%
  \setsansfont{Helvetica Neue}%
}{%
  \IfFontExistsTF{Inter}{\setsansfont{Inter}}{%
    \IfFontExistsTF{Arial}{\setsansfont{Arial}}{}}%
}

% ==== TikZ ====
\usepackage{tikz}
\usetikzlibrary{positioning,arrows.meta,fit,shapes.geometric,calc,%
  decorations.pathmorphing,backgrounds}

\tikzset{
  node-box/.style={draw=cMain,thick,rounded corners=2pt,
    minimum height=8mm,minimum width=18mm,align=center,font=\small},
  node-state/.style={draw=cMain,thick,circle,minimum size=14mm,
    align=center,font=\small},
  node-ok/.style={draw=cOK,thick,fill=cOK!10},
  node-warn/.style={draw=cWarn,thick,fill=cWarn!15},
  node-bad/.style={draw=cBad,thick,fill=cBad!10},
  % 色盲友好的状态机三态：蓝 (正常/Closed) / 橙 (告警/Open) / 灰 (中间/HalfOpen)
  % 避免 red-green 二元区分，改用 hue + lightness 双通道。
  node-state-normal/.style={draw=cMain,thick,fill=cMain!12},
  node-state-alert/.style={draw=cAccent,thick,fill=cAccent!18},
  node-state-transition/.style={draw=cGray,thick,fill=cGray!18},
  arrow/.style={-{Stealth[length=2mm]},thick,draw=cMain},
  arrow-async/.style={-{Stealth[length=2mm]},thick,dashed,draw=cGray},
  arrow-bad/.style={-{Stealth[length=2mm]},thick,draw=cBad},
  arrow-alert/.style={-{Stealth[length=2mm]},thick,draw=cAccent},
  lifeline/.style={dashed,thick,draw=cGray},
}

% ==== 代码 ====
\usepackage{listings}
\lstdefinelanguage{Go}{
  morekeywords={break,case,chan,const,continue,default,defer,else,
    fallthrough,for,func,go,goto,if,import,interface,map,package,range,
    return,select,struct,switch,type,var,nil,true,false,iota},
  morekeywords=[2]{string,int,int64,uint,uint64,bool,byte,error,context},
  sensitive=true,
  morecomment=[l]//,morecomment=[s]{/*}{*/},
  morestring=[b]",morestring=[b]`,
}
\lstdefinelanguage{LuaRedis}{
  morekeywords={and,break,do,else,elseif,end,false,for,function,if,in,
    local,nil,not,or,repeat,return,then,true,until,while},
  morekeywords=[2]{redis,call,KEYS,ARGV,tonumber,HGET,HSET,GET,SET,
    DECR,DECRBY,INCR,INCRBY,EXPIRE,PEXPIRE,ZADD,ZCARD,ZREMRANGEBYSCORE,
    EVAL,EVALSHA},
  sensitive=true,
  morecomment=[l]{--},
  morestring=[b]",morestring=[b]',
}

\lstdefinestyle{base}{
  basicstyle=\ttfamily\scriptsize,
  numbers=left,numberstyle=\tiny\color{cGray},
  keywordstyle=\color{cMain}\bfseries,
  keywordstyle=[2]\color{cAccent},
  stringstyle=\color{cOK},
  commentstyle=\color{cGray}\itshape,
  backgroundcolor=\color{cMute},
  frame=single,framerule=0pt,
  showstringspaces=false,
  breaklines=true,
  tabsize=2,
  columns=fullflexible,
}
\lstset{style=base}

% ==== 三线表与附录页码 ====
\usepackage{booktabs}
\usepackage{appendixnumberbeamer}

% ==== 页脚来源标注 ====
\newcommand{\srcnote}[1]{{\color{cGray}\scriptsize\itshape #1}}

% 关键论点框
\newcommand{\keypoint}[1]{%
  \begin{tcolorbox}[colback=cMute,colframe=cMain,boxrule=0.5pt,arc=2pt]%
  #1\end{tcolorbox}}


\title{库存与防超发}
\subtitle{业务承诺 / 客诉自愈 / 两桶语义 / Saga 退回 / 启动重建}
\author{gomall · 业务 deck 07}
\date{}

\begin{document}

\maketitle

\begin{frame}{今日大纲}
\tableofcontents
\end{frame}

% ============================================================
\section{库存对五个角色意味着什么}
% ============================================================

\begin{frame}{同一个"库存数字"，五个角色五种关心点}
\begin{tabular}{@{}lp{3.0cm}p{5.2cm}@{}}
\toprule
角色 & 关心的事 & 库存系统该回答的问题 \\
\midrule
\textbf{C 端用户}    & 加购了还能不能买       & 提交订单时不能"突然没货" \\
\textbf{商家}        & 我卖了多少 / 还剩多少 & 哪个 SKU 跌破阈值要补货 \\
\textbf{运营 / 平台} & 缺货 SKU 看板 / 补货决策 & 哪些热品长期 reserved 高 \\
\textbf{客服}        & 超卖怎么补偿 / 误锁怎么解释 & 给出可执行话术 + 工单根因定位 \\
\textbf{SRE}         & Redis 跟 DB 一致吗     & 启动重建 / 巡检对账 / 告警分级 \\
\bottomrule
\end{tabular}
\vspace{2mm}
\small "库存"不是一个数字，是\textbf{五条不同业务承诺}的交叉点。本 deck 顺着这五个角色把每件事摊开。
\srcnote{repository/cache/inventory.go:11--14 注释揭示两桶定位}
\end{frame}

\begin{frame}{下单到付款，中间那 15 分钟的业务窗}
\begin{itemize}
  \item 用户点"提交订单"到点"立即支付"，gomall 给\textbf{一段缓冲时间}。
  \item 这段时间里：商品\textbf{不能再被别人买走}，但也\textbf{还没真正卖出}。
  \item 业务侧叫"占着但没消耗"，本仓库用 Redis 两桶建模：available / reserved。
  \item 单桶 \texttt{stock=N} 表达不了这个中间态 —— 要么误锁（用户没付就丢库存），要么超卖（两个人下同一件）。
\end{itemize}
\srcnote{repository/cache/inventory.go:11--14；internal/order/task.go:24 现行口径 15 分钟}
\end{frame}

\begin{frame}{本 deck 七条业务问题}
\begin{enumerate}
  \item 1 个超卖订单的业务后果到底有多严重？工时、补偿、商誉怎么计？
  \item 为什么把库存切成 available + reserved 两桶？业务意义是什么？
  \item Redis 重启 / 副本扩容时，库存数字怎么从 DB 重建？期间能下单吗？
  \item 下单成功但 DB 落库失败，预占的库存会不会泄漏成"凭空缺货"？
  \item 关单（超时 / 用户主动）触发的库存退回，业务上几分钟是合适的？
  \item 客服收到"加购后还有 5 件，下单时说缺货"工单，怎么自查 + 怎么回？
  \item 库存对外业务承诺（SLO）是什么？哪些场景超过 SLO 就要赔？
\end{enumerate}
\srcnote{对应代码 repository/cache/inventory.go + service/inventory/syncer.go + internal/order/service.go + internal/order/cancel.go}
\end{frame}

% ============================================================
\section{超卖：一笔订单的真实业务代价}
% ============================================================

\begin{frame}{1 笔超卖订单 = 一连串的钱、工时、品牌损失}
\begin{tabular}{@{}lp{4.6cm}r@{}}
\toprule
环节 & 业务后果 & 估算成本 \\
\midrule
客服工单         & 用户怒打电话 / 在线投诉，平均 1 单 25 min & 客服工时 \\
用户补偿         & 退款 + 道歉券（10 / 30 / 50 元档）        & 10--50 元 \\
商家发不出货     & 商家 SLA 违约罚款 / 自掏腰包补货          & 商品成本 + 罚金 \\
平台代偿         & 同档替代品差价 / 自营仓兜底              & 差价 + 物流 \\
品牌口碑         & 社交平台截图传播，新用户转化下滑           & 长期 GMV 损失 \\
\midrule
\textbf{单笔 100 元商品超卖} & 直接现金 \textbf{≈ 60--150 元} + 不可量化口碑 & 远超商品本身利润 \\
\bottomrule
\end{tabular}
\vspace{2mm}
\small 这是为什么库存系统的工程目标是\textbf{宁可误锁不可超卖} —— 误锁是体验事件（用户去别家买），超卖是法律 / 合同 / 财务 / 品牌四件事一起爆。
\srcnote{repository/cache/inventory.go:11--14 注释直接讲两桶模型的业务红线}
\end{frame}

\begin{frame}{零超卖怎么证明：500 个用户同时抢 100 件}
\begin{tabular}{@{}lr@{}}
\toprule
业务指标 & 实测 \\
\midrule
活动准备库存 & 100 件 \\
同时下单用户数 & 500 \\
\textbf{成功下单数} & \textbf{恰好 100} \\
拒单数（业务码：库存不足） & 400 \\
活动结束 available 桶 & 0 \\
活动结束 reserved 桶 & 100（等待支付） \\
\bottomrule
\end{tabular}
\vspace{2mm}
\small 业务侧解读：\textbf{"恰好 100"是写死的 CI 红线}（少 1 个少卖、多 1 个超卖都不允许）。这条断言代表的不是技术指标，而是\textbf{对商家"我备货 100 件就只卖 100 件"的承诺}。
\srcnote{repository/cache/inventory\_test.go:52--82 NoOversellUnderConcurrency；stressTest/REPORT.md \S 6}
\end{frame}

\begin{frame}{尾延迟为什么是业务问题：Redis Lua vs DB 行锁}
\begin{tabular}{@{}lrr@{}}
\toprule
& Redis Lua & DB \texttt{SELECT FOR UPDATE} \\
\midrule
平均延迟 & 1.24 ms & 1.29 ms \\
p95 延迟 & 3.52 ms & 3.65 ms \\
\textbf{最差用户体验} & \textbf{136 ms} & \textbf{453 ms} \\
吞吐 (RPS) & 51{,}362 & 50{,}142 \\
\bottomrule
\end{tabular}
\vspace{2mm}
\small 平均 / p95 看不出差距，但\textbf{最差那个用户}从 136ms 涨到 453ms —— 这是 3.3 倍体验差。\\
业务后果：大促 100k 并发时，DB 锁路径会让\textbf{尾部几百个用户}看到 1--2 秒转圈，转化率掉点；Redis Lua 路径把"最差体验"压在 150ms 以内。
\srcnote{stressTest/REPORT.md \S 6 优惠券 Redis Lua vs DB 锁实测}
\end{frame}

% ============================================================
\section{两桶库存的业务语义}
% ============================================================

\begin{frame}{available / reserved 翻译成业务话}
\begin{tikzpicture}[scale=1]
  \draw[thick,draw=cMain,fill=cMain!15] (0,0) rectangle (6,1);
  \node[font=\small,anchor=west] at (0.2,0.5) {\textbf{available} = 80};
  \node[font=\scriptsize\itshape,anchor=east,color=cGray] at (5.8,0.5) {"还能卖给下一个人"};

  \draw[thick,draw=cAccent,fill=cAccent!18] (0,-1.4) rectangle (1.5,-0.4);
  \node[font=\small,anchor=west] at (0.2,-0.9) {\textbf{reserved} = 20};
  \node[font=\scriptsize\itshape,anchor=west,color=cGray] at (2,-0.9) {"答应给某人但还没收钱"};

  \draw[thick,draw=cGray,fill=cMute] (0,-2.8) rectangle (6,-1.8);
  \node[font=\small,anchor=west] at (0.2,-2.3) {\textbf{product.num} = 100（商家备货总量）};

  \draw[arrow] (4,-0.1) -- node[right,font=\scriptsize]{下单}  (1,-0.3);
  \draw[arrow,draw=cOK] (1,-1.5) -- node[left,font=\scriptsize]{超时未付 / 取消}  (4,-1.7);
\end{tikzpicture}
\vspace{2mm}
\small \textbf{客服侧的语义价值}：当用户问"为什么我加购后看见还有 5 件，下单却说没货"，客服可以直接答\\
\quad "其他 5 件被别的客户下单未付款占着，30 分钟内若对方不付款会自动放回来"。\\
两桶不是技术名词，是\textbf{对客服可解释的业务模型}。
\srcnote{repository/cache/inventory.go:16--22 两个 key 命名规范}
\end{frame}

\begin{frame}{物理库存 vs 可售库存：商家 / 客服 / 财务的语言对齐}
\footnotesize
\begin{tabular}{@{}llp{6.8cm}@{}}
\toprule
桶 & 谁在看 & 业务含义 \\
\midrule
\textbf{available}  & C 端 / 商详页   & "还能下单的数字" —— 唯一可对外曝光 \\
\textbf{reserved}   & 客服 / 商家后台 & "已下单未支付" —— 客服查"为什么缺货" \\
committed           & 财务 / 商家     & "已支付未发货" —— 进 GMV / 应发货量 \\
in\_transit         & 物流 (WMS)      & 在途，本仓库未实现 \\
defective           & 仓储 (WMS)      & 残次，不入可售 \\
\bottomrule
\end{tabular}
\smallskip
\keypoint{\footnotesize \textbf{业务对齐三原则}：(1) C 端\textbf{只看 available} —— 商详页绝不显示 reserved（会让用户"焦虑性下单"）。(2) 商家后台看 available + reserved —— 知道"承诺出去了多少"才好备货。(3) 财务对账盯恒等式：available + reserved + committed = product.num。}
\srcnote{repository/cache/inventory.go:127--138 GetStockSnapshot 给巡检 / 财务用}
\end{frame}

\begin{frame}{三态生命周期：业务侧四件事}
\begin{tikzpicture}[node distance=8mm and 16mm]
  \node[node-state,node-state-normal] (av)  {available};
  \node[node-state,node-state-transition,right=of av] (rv) {reserved};
  \node[node-state,node-state-alert,right=of rv] (cm) {committed};
  \node[node-state,node-state-alert,below=of cm] (it) {in\_transit};
  \node[node-state,node-state-normal,below=of av] (df) {defective};
  \draw[arrow] (av) -- node[above,font=\tiny]{下单 reserve} (rv);
  \draw[arrow] (rv) -- node[above,font=\tiny]{支付 commit} (cm);
  \draw[arrow,draw=cOK,bend left=35] (rv.north) to node[above,font=\tiny\color{cOK}]{超时 / 取消 release} (av.north);
  \draw[arrow] (cm) -- node[right,font=\tiny]{发货 ship} (it);
  \draw[arrow-bad,bend left=30] (it.west) to node[above,font=\tiny\color{cBad}]{退货回仓} (av.east);
  \draw[arrow-bad] (it) -- node[below,font=\tiny\color{cBad}]{破损} (df);
\end{tikzpicture}
\vspace{1mm}
\begin{itemize}\setlength\itemsep{-1mm}
  \item \textbf{下单 reserve}：占住一份库存，承诺给该用户。
  \item \textbf{支付 commit}：钱到账，库存真正消耗，不再可退回 available。
  \item \textbf{超时 / 取消 release}：承诺取消，库存退回再卖给下一个人。
  \item \textbf{发货 ship / 退货 / 破损}：归 WMS 仓储系统，本仓库\textbf{不做}（路线图）。
\end{itemize}
\srcnote{repository/cache/inventory.go:29--66 三段 Lua 脚本完整定义}
\end{frame}

% ============================================================
\section{下单 reserve：业务承诺的工程化}
% ============================================================

\begin{frame}[fragile]{8 行 Lua = "宁可误锁不可超卖"的工程化}
\begin{lstlisting}[language=LuaRedis,firstnumber=32,
  caption={\srcnote{repository/cache/inventory.go:32--44 reserveScript}}]
local avail = redis.call('GET', KEYS[1])    -- 1. 读 available
if avail == false then
    return -2                                -- 2. 库存未初始化 -> ErrStockNotInit
end
local need = tonumber(ARGV[1])
if tonumber(avail) < need then
    return -1                                -- 3. 不足 -> ErrStockInsufficient
end
redis.call('DECRBY', KEYS[1], need)          -- 4. available -= n
redis.call('INCRBY', KEYS[2], need)          -- 5. reserved  += n
return 1                                     -- 6. 占位成功
\end{lstlisting}
\small 业务侧关键：第 2 / 第 3 步\textbf{在扣减之前}，宁可拒单也不允许扣穿 —— 这就是"宁可误锁"的工程化。
\textbf{诚实标注}：当前代码里 \texttt{ErrStockInsufficient} 是 Go 错误对象，\textbf{未映射到独立业务码}（不是 50001，50001 是 OSS 错误），上层 handler 返回的是通用 ERROR；建议路线图加 \texttt{ErrStockInsufficient = 90001}（80001--80003 已被满减引擎业务码占用）。
\srcnote{pkg/e/code.go 当前未定义库存业务码；repository/cache/inventory.go:25 ErrStockInsufficient 错误对象}
\end{frame}

\begin{frame}[fragile]{下单链路 + 失败自愈：业务侧三步走}
\small 链路：\texttt{用户点提交 $\to$ Redis 占库存 $\to$ DB 写订单 + outbox $\to$ MQ 排关单延迟}。\\
任一步失败：\textbf{已占的库存必须退回}，否则商家会"凭空缺货"。
\begin{lstlisting}[language=Go,firstnumber=101,
  caption={\srcnote{internal/order/service.go:101--155 reserve + tx + 补偿 release（满减结算段略）}}]
// 1) Redis 预扣: available -> reserved，失败直接拒单
if err = cache.ReserveStock(ctx, req.ProductID, int64(req.Num)); err != nil {
    return nil, err                       // ErrStockInsufficient
}
// 2) 同事务写订单 + 应用满减 + 写 outbox（ApplyDiscountInTx 段略）
err = dao.NewDBClient(ctx).Transaction(func(tx *gorm.DB) error {
    if e := NewOrderDaoByDB(tx).CreateOrder(order); e != nil { return e }
    return outbox.NewOutboxDaoByDB(tx).Insert(/* OrderCreated */)
})
if err != nil {
    // 3) 事务失败：必须释放刚扣的预占（Saga 补偿）
    if relErr := cache.ReleaseReservation(ctx,
            req.ProductID, int64(req.Num)); relErr != nil {
        util.LogrusObj.Errorf("release on tx failure: %v", relErr)
    }
    return
}
\end{lstlisting}
\small \textbf{先 Redis 后 DB} 的业务理由：DB 慢的时候 Redis 直接拒单，不把过载请求压给 DB；DB 是真理之源，热写打进 DB 会让"全平台"一起慢。
\end{frame}

% ============================================================
\section{Saga 退回：避免"凭空缺货"事故}
% ============================================================

\begin{frame}{业务事故：DB 落库失败而预占未退 = 库存泄漏}
\begin{center}
\begin{tikzpicture}[node distance=8mm and 14mm]
  \node[node-box,node-warn] (s1) {step1: Redis 占位\\available--, reserved++};
  \node[node-box,right=of s1] (s2) {step2: DB 写订单\\order + outbox};
  \node[node-box,right=of s2] (s3) {step3: MQ 排关单};
  \node[node-box,below=of s2,node-bad] (c1) {补偿: 退回\\reserved--, available++};
  \draw[arrow] (s1) -- (s2);
  \draw[arrow] (s2) -- (s3);
  \draw[arrow-bad] (s2) -- node[right,font=\tiny\color{cBad}]{tx 失败} (c1);
  \draw[arrow-bad,bend right=20] (c1.west) to node[below,font=\tiny\color{cBad}]{回滚 step1}(s1.south);
\end{tikzpicture}
\end{center}
\vspace{1mm}
\small \textbf{如果不补偿}：available 数字永远比真实可售小，长期堆积\textbf{所有 Redis 库存被锁死} —— 商家明明备了 100 件，前台只显示"还剩 0" —— 这是真实 P1 事故。\\
客服话术：用户来投诉时，先查 outbox 是否有同 orderNum 的 \texttt{OrderCreated} 事件 + Redis snapshot 的 reserved 是否偏高，判断是否触发补偿。
\srcnote{internal/order/service.go:148--154 显式调用 ReleaseReservation 补偿}
\end{frame}

\begin{frame}[fragile]{commit / release 成对设计的业务理由}
\begin{lstlisting}[language=LuaRedis,firstnumber=48,
  caption={\srcnote{repository/cache/inventory.go:48--66 commit + release}}]
-- commitScript：支付成功，reserved 真正消耗（钱到账，不退）
local r = redis.call('GET', KEYS[1])           -- KEYS[1]=reserved
if r == false or tonumber(r) < tonumber(ARGV[1]) then
    return -1
end
redis.call('DECRBY', KEYS[1], ARGV[1])         -- reserved -= n
return 1
-- releaseScript：取消 / 超时，reserved 退回 available
local r = redis.call('GET', KEYS[1])           -- KEYS[1]=reserved
if r == false or tonumber(r) < tonumber(ARGV[1]) then
    return -1
end
redis.call('DECRBY', KEYS[1], ARGV[1])         -- reserved -= n
redis.call('INCRBY', KEYS[2], ARGV[1])         -- available += n
return 1
\end{lstlisting}
\small \textbf{commit 只动 reserved}（钱付了，承诺兑现，库存不再属于"可退回"池）；\textbf{release 两桶都动}（承诺取消，库存重新可卖）。check $r \geq n$ 避免 reserved 跌负 —— 跌负 = 同一笔订单退了两次 = 库存凭空多出。
\end{frame}

% ============================================================
\section{关单释放：15 / 30 / 60 分钟的业务权衡}
% ============================================================

\begin{frame}{关单超时定多少分钟：商家利用率 vs 用户犹豫时间}
\begin{tabular}{@{}lp{4cm}p{4cm}@{}}
\toprule
档位 & 用户体验 & 商家库存利用率 \\
\midrule
\textbf{15 分钟} & 大促时用户来不及付款 / 充值，关单率偏高 & 库存周转最快，热品多卖一轮 \\
\textbf{30 分钟}（业界中位） & 给用户去切支付方式 / 看评论的时间 & 周转中等 \\
\textbf{60 分钟}            & 用户犹豫期长，关单率低                & 热品被锁太久，错过下一波流量 \\
\bottomrule
\end{tabular}
\vspace{2mm}
\small gomall 当前现状：\\
\quad - RMQ TTL 延迟队列 \texttt{OrderCancelDelay = 30 \* time.Minute} \srcnote{repository/rabbitmq/order\_delay.go:22}\\
\quad - Cron 兜底扫表 \texttt{GetTimeoutOrders(15, 100)} 视 15min 为超时 \srcnote{internal/order/task.go:24}\\
\textbf{这两处口径不一致}（README 已自爆）：MQ 30 分钟到点关、Cron 15 分钟见就关，实际生效以\textbf{先到的为准 = 15 分钟}。修法是统一到 \texttt{consts.OrderCancelTimeout}，避免商家投诉"为什么我设了 30 分钟实际 15 分钟就关了"。
\srcnote{consts 路线图：单一来源 OrderCancelTimeout 常量}
\end{frame}

\begin{frame}[fragile]{关单释放的幂等保证：客服可以放心"再点一次"}
\small 业务背景：客服在工单系统点"立即关单"按钮，可能因为前端超时连点两次；RMQ 重投也会同 orderNum 触发两次；Cron 兜底也会撞上 —— 必须\textbf{多次执行结果一致}，库存不能退第二次。
\begin{lstlisting}[language=Go,firstnumber=31,
  caption={\srcnote{internal/order/cancel.go:31--63 关单 + 库存退回}}]
err = baseDao.DB.Transaction(func(tx *gorm.DB) error {
    ok, err := NewOrderDaoByDB(tx).CloseOrderWithCheck(orderNum)
    if err != nil { return err }
    if !ok { return nil }                  // 1. 已关 -> 短路 no-op
    closed = true
    return outbox.NewOutboxDaoByDB(tx).Insert(/* OrderCancelled 载荷含 promo 字段 */)
})
if err != nil { return err }
if !closed { return nil }                  // 2. 没真关 -> 不退库存

if relErr := cache.ReleaseReservation(ctx,
    order.ProductID, int64(order.Num)); relErr != nil {
    util.LogrusObj.Errorf("release on cancel failed: %v", relErr)
}
\end{lstlisting}
\small \textbf{closed 标志位是幂等的关键}：\texttt{CloseOrderWithCheck} 只对 WaitPay（待付款）生效，二次调用 ok=false 整段 no-op，库存\textbf{不会退第二次}。客服 / RMQ / Cron / 用户主动取消 4 个调用口，统一从这个入口走。
\end{frame}

\begin{frame}{关单延迟链路：MQ 主动 + Cron 兜底}
\begin{tikzpicture}[node distance=8mm and 14mm]
  \node[node-box] (ord) {下单};
  \node[node-box,right=of ord,node-warn] (mq) {RMQ TTL\\30min};
  \node[node-box,right=of mq] (worker) {CancelWorker};
  \node[node-box,below=of mq] (cron) {Cron 兜底\\15min 扫表};
  \node[node-box,right=of worker,node-ok] (rel) {ReleaseReservation\\reserved -> available};
  \draw[arrow] (ord) -- (mq);
  \draw[arrow] (mq) -- node[above,font=\tiny]{到点投递} (worker);
  \draw[arrow] (cron) -- (worker);
  \draw[arrow] (worker) -- (rel);
\end{tikzpicture}
\vspace{2mm}
\small \textbf{业务侧为什么双保险}：\\
- 单 RMQ：进程重启 / 消息丢 = 漏关 = 商家"凭空缺货"。\\
- 单 Cron：5min 扫一次，最差用户付完款 5min 才看到 available 回补。\\
- 双保险：MQ 主路径秒级，Cron 兜底分钟级，两路汇到同一个幂等 \texttt{CancelUnpaidOrder} 入口。
\srcnote{repository/rabbitmq/order\_delay.go:22 OrderCancelDelay；internal/order/task.go:24 Cron 15min}
\end{frame}

% ============================================================
\section{启动重建：服务重启时的业务窗口}
% ============================================================

\begin{frame}{业务场景：Redis 重启，库存数字从哪儿来？}
\begin{tikzpicture}[node distance=8mm and 16mm]
  \node[node-box] (boot) {gomall 启动};
  \node[node-box,right=of boot] (seed) {SeedFromDB};
  \node[node-box,above right=of seed,node-warn] (db) {MySQL\\product.num\\（真理之源）};
  \node[node-box,below right=of seed,node-ok] (rd) {Redis\\stock:available:*};
  \draw[arrow] (boot) -- (seed);
  \draw[arrow] (db) -- node[above,font=\tiny]{游标分页扫表} (seed);
  \draw[arrow] (seed) -- node[below,font=\tiny]{InitStock（跳过已存在）} (rd);
\end{tikzpicture}
\vspace{2mm}
\begin{itemize}
  \item \textbf{真理之源}是 MySQL 的 \texttt{product.num}（商家后台维护，binlog 持久化）。
  \item Redis 是\textbf{运行时账本}，重启会丢，必须能从 DB 复制回来。
  \item \textbf{已经在跑的实例不能被覆盖}（多副本滚动重启时，新实例不能抹掉旧实例的 reserved）。
\end{itemize}
\srcnote{initialize/inventory.go:11--17；service/inventory/syncer.go:14}
\end{frame}

\begin{frame}[fragile]{SeedFromDB 主循环：游标分页 + EXISTS 跳过}
\begin{lstlisting}[language=Go,firstnumber=14,
  caption={\srcnote{service/inventory/syncer.go:14--48 SeedFromDB}}]
func SeedFromDB(ctx context.Context, batchSize int) error {
    if batchSize <= 0 { batchSize = 500 }            // 默认页 500
    d := dao.NewDBClient(ctx); var lastID uint
    for {
        var rows []*product.Product
        q := d.Model(&product.Product{}).Order("id ASC").Limit(batchSize)
        if lastID > 0 { q = q.Where("id > ?", lastID) }  // 游标分页
        if err := q.Find(&rows).Error; err != nil { return err }
        if len(rows) == 0 { break }
        for _, p := range rows {
            lastID = p.ID
            exists, _ := cache.RedisClient.Exists(ctx,
                cache.StockAvailableKey(p.ID)).Result()
            if exists == 1 { continue }              // 跳过运行期数据
            cache.InitStock(ctx, p.ID, int64(p.Num)) // 覆盖写 0
        }
    }
    return nil
}
\end{lstlisting}
\small \textbf{两个业务保护}：(1) 游标分页避免重启时把 DB CPU 打满；(2) EXISTS 跳过保护"运行期已经在卖的实例"，防止新副本起来抹平 reserved。
\end{frame}

\begin{frame}{重建窗口里能不能下单？客服怎么解释？}
\begin{tabular}{@{}lp{4.0cm}p{4cm}@{}}
\toprule
窗口阶段 & C 端体感 & 客服话术 \\
\midrule
0--100s（重建中） & 部分商品返"库存未初始化" & "系统刚刚重启，1 分钟内会恢复" \\
重建完成 & 全量 available 已回灌 & 正常下单 \\
重建失败（DB 慢） & 持续返"库存未初始化" & "高峰期个别商品暂时不可下单，建议稍后" \\
\bottomrule
\end{tabular}
\vspace{2mm}
\small \textbf{业务取舍}：1M SKU 重建 ≈ 100s。期间下单返 \texttt{ErrStockNotInit} 触发降级 banner —— \textbf{宁可少卖 1 分钟也不允许超卖}。\\
\textbf{SRE 预案}：大促前一天手动跑 SeedFromDB 灌全量；活动前 5min 锁商家改库存入口避免双写竞态。
\srcnote{repository/cache/inventory.go:69--75 InitStock；syncer.go:14 SeedFromDB}
\end{frame}

% ============================================================
\section{客服 SOP：三类高频工单的自查路径}
% ============================================================

\begin{frame}{工单 A："我加购后还有 5 件，下单时说缺货"}
\textbf{业务真相}：商详页缓存了 5min 前的快照，实际 available 已经被别人下单变成 0。
\vspace{1mm}
\begin{itemize}
  \item \textbf{客服自查 3 步}：
  \begin{enumerate}
    \item 查商品 ID $\to$ \texttt{GetStockSnapshot} 看 available / reserved；
    \item 若 reserved 很高（如 80\%）$\to$ 大量预占未付款，30 分钟内会自愈；
    \item 若 available = 0 reserved = 0 $\to$ 真的卖完了，建议同档替代。
  \end{enumerate}
  \item \textbf{话术}：\\
  \quad - reserved 偏高：\textit{"该商品 80\% 库存被其他用户下单未付款占着，半小时内若对方不付款会回补，建议您 30 分钟后再试。"}\\
  \quad - 真没了：\textit{"该商品已全部售出，为您推荐同款……"}
  \item \textbf{自愈机制}：MQ TTL + Cron 双保险确保 30 分钟内一定释放，客服不需要人工介入。
\end{itemize}
\srcnote{repository/cache/inventory.go:127--138 GetStockSnapshot；internal/order/cancel.go:61 release 调用}
\end{frame}

\begin{frame}{工单 B："我下单成功了为什么显示缺货"}
\textbf{业务真相}：该用户的订单成功进了 reserved，但商详页 available = 0，用户在订单列表看到"已下单"、在商详页看到"已售罄"，体感矛盾。
\vspace{1mm}
\begin{itemize}
  \item \textbf{客服自查}：
  \begin{enumerate}
    \item 查 orderNum $\to$ 订单状态 = WaitPay；
    \item 查 \texttt{GetStockSnapshot} reserved 包含该用户的 num；
    \item \textbf{确认是正常占位}，不是事故。
  \end{enumerate}
  \item \textbf{话术}：\textit{"您的订单已成功提交（在 reserved 状态），库存已经为您锁定。商详页显示售罄表示其他用户暂时不能再下单。请 30 分钟内完成支付。"}
  \item \textbf{业务设计意图}：商详页\textbf{只看 available} 是有意的 —— 不能给"焦虑性下单"开口子。
\end{itemize}
\srcnote{repository/cache/inventory.go:11--14 注释指明 available 是对外可售}
\end{frame}

\begin{frame}{工单 C："我退款了为什么库存没还回去"}
\textbf{业务真相}：用户在支付后退款，库存应当从 committed 回到 available 或进入 defective。退款状态机已落地（\texttt{internal/refund/service.go}，含 \texttt{order.refunded} 事件），但\textbf{库存回流仍未实现}（路线图阶段）。
\vspace{1mm}
\begin{itemize}
  \item \textbf{现行处理}：人工操作 —— SRE 用 \texttt{InitStock} 人工覆盖 available（误锁可接受、超卖不可）。
  \item \textbf{话术}：\textit{"您的退款已生效，库存回填会在 24 小时内由商家手动确认完成（避免在途货退回错位）。"}
  \item \textbf{路线图}：\\
  \quad - \texttt{order.refunded} outbox 事件已落地 $\to$ 缺消费侧触发 \texttt{ReleaseFromCommitted} Lua（commit 的逆操作）。\\
  \quad - 退货寄回 / 残次品分流 $\to$ WMS 仓储系统对接。\\
  \quad - 售后工单 $\to$ 评价 $\to$ 二次售后的全链路在 deck 09 路线图章节。
\end{itemize}
\srcnote{退款状态机已实现 internal/refund/service.go:94--135 ApproveRefund；库存回流扩展点：ReleaseFromCommitted Lua}
\end{frame}

\begin{frame}{业务码 / 错误的客服对照表（含诚实标注）}
\footnotesize
\begin{center}
\begin{tabular}{@{}p{4.0cm}p{3.6cm}p{3.4cm}@{}}
\toprule
状态码 / 错误 & 现行实现 & 客服话术 \\
\midrule
\textbf{ErrStockInsufficient}（库存不足） & Go 错误对象，\textbf{未映射独立业务码} & "本商品已售罄" \\
\textbf{ErrStockNotInit}（库存未初始化） & 重建窗口内出现 & "系统恢复中请稍候" \\
60002 ErrIdempotencyInProgress & 同 key 重复下单 & "您已下单成功" \\
70001 ErrRateLimitExceeded & 限流命中 & "活动太火爆请稍后" \\
70002 ErrCircuitOpen      & 支付熔断 Open & "支付服务暂忙 1 分钟后" \\
\bottomrule
\end{tabular}
\end{center}
\smallskip
\keypoint{\footnotesize \textbf{诚实标注}：50001 在本仓库是 \texttt{ErrorOss}（对象存储错误），\textbf{不是}库存不足；库存不足现仅 Go 错误对象 \texttt{cache.ErrStockInsufficient}，handler 返回通用 ERROR。\textbf{路线图}：给库存独立业务码 \texttt{ErrStockInsufficient = 90001} + \texttt{ErrStockNotInit = 90002}（80001--80003 已被满减引擎占用），客服 / 监控 / 链路追踪都能精准区分。}
\srcnote{pkg/e/code.go：50001=ErrorOss / 70001 限流 / 70002 熔断；repository/cache/inventory.go:25 ErrStockInsufficient 仅 Go 错误对象}
\end{frame}

% ============================================================
\section{业务承诺 SLO 与超卖应急}
% ============================================================

\begin{frame}{库存子系统的业务 SLO}
\begin{tabular}{@{}lll@{}}
\toprule
业务承诺 & 指标 & 目标 \\
\midrule
\textbf{零超卖} & 单测 NoOversell 必须绿 & 100\%（CI 红线） \\
下单库存判断 p99 & ReserveStock 端到端 & < 10 ms \\
最差体验 & Redis Lua max & < 200 ms（实测 136 ms） \\
Saga 退回窗口 & DB 失败到 release 完成 & < 1s \\
启动重建 SLA & 1M SKU 全量重建 & < 120s \\
关单释放及时性 & 30min TTL + 15min Cron 兜底 & MQ p99 < 31min \\
对账偏差告警 & available + reserved + committed = product.num & 偏差 > 1\% 告警 \\
\bottomrule
\end{tabular}
\vspace{2mm}
\small "零超卖"在工程上 = \texttt{TestInventory\_NoOversellUnderConcurrency} 写死 \texttt{success == exactly 100}；少卖 / 超卖 CI 都红 —— \textbf{业务红线工程化}。
\srcnote{repository/cache/inventory\_test.go:52--82；stressTest/REPORT.md \S 6}
\end{frame}

\begin{frame}{发生超卖了：4 小时应急 SOP}
\footnotesize
\begin{tabular}{@{}lp{6.4cm}p{3.6cm}@{}}
\toprule
时间窗 & 动作 & 责任方 \\
\midrule
T+0 min   & 监控告警：可发货量 < 已支付订单量 & SRE on-call \\
T+5 min   & 锁定问题 SKU：停下单入口（动态开关） & SRE 改路由 \\
T+15 min  & 跑对账脚本：超卖量 + 影响订单数 & DBA + 库存 owner \\
T+30 min  & 客服话术下发：哪些订单发货 / 赔付 & 客服中心 \\
T+1h      & 给被超卖用户发\textbf{补偿券 + 致歉短信}  & 营销 + 客服 \\
T+4h      & 复盘 + 修复脚本上线 & 库存 owner \\
\bottomrule
\end{tabular}
\smallskip
\keypoint{\footnotesize \textbf{关键指标}：T+5 min 内能锁单是底线。商详页"下单入口"必须是动态开关，不能等发版（动态开关本身也是路线图）。\textbf{补偿档位}：< 24h = 退款 + 10 元券；24--72h = 退款 + 30 元券 + 致歉短信；> 72h = 退款 + 50 元券 + 同档替代 + 1v1 客服。}
\srcnote{应急 SOP；对账基于 GetStockSnapshot + DB SUM 比对}
\end{frame}

\begin{frame}{对账恒等式 + 巡检阈值}
\textbf{恒等式}：\texttt{available + reserved + committed = product.num}\\
前两项来自 \texttt{GetStockSnapshot}，committed = DB \texttt{SUM(num)}（type 取已付各态，自 OrderWaitShip 起）。
\vspace{1mm}
\begin{tabular}{@{}lp{6cm}@{}}
\toprule
偏差 & 处理 \\
\midrule
< 0.1\% & 日志记录，不告警 \\
0.1\%--1\% & 钉钉提醒，T+1 复查 \\
1\%--5\%  & 钉钉 P3，立即排查 \\
\textbf{> 5\%} & \textbf{电话 P1}，启动超卖应急 SOP \\
大促前 > 0.1\% & 拦截活动开始 \\
\bottomrule
\end{tabular}
\vspace{2mm}
\small \textbf{校验时机}：10min 抽样 100 个热门 SKU + 凌晨 3:00 全量对账。\\
\textbf{修复}：SRE 用 \texttt{InitStock(pid, 真实可售)} 人工覆盖；优先选"误锁"侧而非"超卖"侧 —— 误锁可挽回，超卖只能补偿。
\srcnote{repository/cache/inventory.go:127--138 GetStockSnapshot}
\end{frame}

% ============================================================
\section{业务边界与路线图}
% ============================================================

\begin{frame}{本库存系统不做什么（诚实清单）}
\begin{itemize}
  \item \textbf{不做多仓库}：当前 Redis key 是 \texttt{stock:available:\{pid\}}，没有 warehouse 维度；多仓分布 + 智能选仓 + 跨仓凑单是路线图。
  \item \textbf{不做 SKU 多规格}：当前以 product\_id 代 SKU，颜色 / 尺寸 / 套餐没有独立扣减；扩展需新增 sku\_id + 订单表迁移。
  \item \textbf{不做期货 / 排队购}：普通商品库存必须实物到仓才上架（预售定金 / 尾款两段式已另行落地，见 deck 15）。
  \item \textbf{不做跨仓调拨 / 在途回流}：发货后归 WMS（本仓库未实现），退货寄回 + 残次品分流缺失。
  \item \textbf{不做商家自助补货 API}：商家想加 100 件库存只能找运营手动改 \texttt{product.num} + 跑 \texttt{InitStock}。
  \item \textbf{不做冻结库存}：质押 / 财务冻结 / 法务冻结这些 B 端需求都没建模。
  \item \textbf{不做防黄牛单用户上限}：reserveScript 暂未加 owned $<$ perUser check（repository/cache/coupon.go:30--44 claimScript 已实现该范式可复用）。
\end{itemize}
\srcnote{docs/architecture/feature-matrix.md 路线图全表}
\end{frame}

\begin{frame}{业务路线图：库存能力扩展}
\begin{tabular}{@{}lp{5.2cm}l@{}}
\toprule
能力 & 业务价值 & 优先级 \\
\midrule
独立库存业务码 90001/90002 & 客服 / 监控 / 链路精准分类 & P0 \\
库存预警事件 & SKU < 10 件推送商家 & P1 \\
商家自助补货 API & 商家后台直接加库存 & P1 \\
运营缺货 SKU 看板 & 补货决策可视化 & P1 \\
售后库存回流 & 退款 $\to$ committed 退回 available & P1 \\
SKU 多规格扩展 & 颜色 / 尺寸 / 套餐独立扣减 & P2 \\
多仓库 + 智能选仓 & 就近发货 + 凑单优化 & P2 \\
单用户限购 Lua & 防黄牛批量下单 & P2 \\
秒杀独立 Redis & 大促隔离主集群 & P2 \\
冷热库存分离 & 普通商品 / 秒杀物理隔离 & P3 \\
\bottomrule
\end{tabular}
\vspace{2mm}
\small \textbf{P0 优先级}：把库存错误从 Go 错误对象升级到独立业务码 —— 不写代码也能在 5min 内做完。\\
\textbf{P1 三件}（库存预警 / 自助补货 / 缺货看板）是"商家入驻"业务的核心 —— 直接影响 GMV。
\srcnote{repository/cache/coupon.go:30--44 已实现的 owned check 可复用到 reserveScript}
\end{frame}

\begin{frame}{业务承诺：超卖 vs 误锁，两条不同的红线}
\begin{tabular}{@{}lll@{}}
\toprule
事故 & 用户体感 & 业务代价 \\
\midrule
\textbf{超卖}     & 下单成功，无法发货  & 真金白银赔 + 商誉 \\
\textbf{误锁}     & 显示"已售罄"        & 转化率轻微下降 \\
\bottomrule
\end{tabular}
\vspace{3mm}
\begin{itemize}
  \item 超卖 = 法律 / 合同事件，\textbf{零容忍}（NoOversell 单测是这条红线的工程化）。
  \item 误锁 = 体验事件，\textbf{可以容忍}（高峰过后 / 关单超时后自愈）。
  \item 这条不对称是\textbf{所有设计选择的源头}：选 Redis 单线程、选先 reserve 后 commit、选 EXISTS 跳过、选先 Redis 后 DB —— 全部都在向"宁可误锁不可超卖"这一侧倾斜。
\end{itemize}
\srcnote{repository/cache/inventory.go:11--14 注释揭示这条业务红线}
\end{frame}

% ============================================================
\appendix
% ============================================================

\begin{frame}{Q\&A（上）：业务承诺 / 工单 / 应急}
\textbf{Q1}：1 笔超卖订单的真实业务代价多大？\\
\hspace{4mm}\small A：100 元商品超卖 ≈ 60--150 元直接现金（客服工时 25min + 补偿 10--50 元 + 商家差价 / 罚金 + 平台代偿）+ 不可量化的品牌口碑。这是为什么"零超卖"是 CI 红线。

\vspace{1mm}
\textbf{Q2}：用户工单"加购后还有 5 件下单却说缺货"，客服怎么自查？\\
\hspace{4mm}\small A：3 步 —— (1) \texttt{GetStockSnapshot} 看 available / reserved；(2) reserved 高 = 别人占着未付款，30min 自愈；(3) 都为 0 = 真卖完，推同款。

\vspace{1mm}
\textbf{Q3}：用户工单"下单成功为什么显示缺货"，是不是事故？\\
\hspace{4mm}\small A：不是。该用户已进 reserved，商详只看 available，是\textbf{有意的产品设计}（避免焦虑性下单）。话术：库存已为您锁定，30 分钟内付款。

\vspace{1mm}
\textbf{Q4}：发生超卖了应急多久？\\
\hspace{4mm}\small A：4 小时 SOP。底线：T+5min 锁定问题 SKU（动态关下单入口），T+30min 发客服话术，T+1h 发补偿券。最关键是"下单入口动态可关"，不能等发版。
\end{frame}

\begin{frame}[allowframebreaks]{Q\&A（下）：业务边界 / 路线图 / 诚实标注}
\footnotesize
\textbf{Q5}：50001 是不是缺货码？\\
\hspace{4mm} A：\textbf{不是}。50001 在本仓库是 \texttt{ErrorOss}（对象存储错误）。库存不足现在只是 \texttt{cache.ErrStockInsufficient} Go 错误对象，未映射独立业务码。路线图：升 90001（80001--80003 已被满减引擎占用）。代码：\texttt{pkg/e/code.go} + \texttt{repository/cache/inventory.go:25}。

\vspace{1mm}
\textbf{Q6}：当前关单超时是 15 还是 30 分钟？\\
\hspace{4mm} A：\textbf{口径不一致}（README 自爆）。RMQ TTL 30min，Cron 兜底 15min，实际生效 = min(两者) = 15min。修法：统一到 \texttt{consts.OrderCancelTimeout}。代码：\texttt{repository/rabbitmq/order\_delay.go:22} + \texttt{internal/order/task.go:24}。

\vspace{1mm}
\textbf{Q7}：用户退款后库存怎么回？\\
\hspace{4mm} A：当前\textbf{没有自动回流}（路线图）。SRE 用 \texttt{InitStock} 人工覆盖。\texttt{order.refunded} outbox 事件已落地，消费侧 \texttt{ReleaseFromCommitted} Lua 库存回流是 P1 路线图。

\vspace{1mm}
\textbf{Q8}：商家想加库存怎么操作？\\
\hspace{4mm} A：当前没有商家自助 API，得找运营手动改 \texttt{product.num} + 跑 \texttt{InitStock}。商家自助补货 + 缺货 SKU 看板 + 库存预警事件三件是"商家入驻"P1 路线图。

\vspace{1mm}
\textbf{Q9}：怎么证明 500 并发不会超卖？\\
\hspace{4mm} A：\texttt{NoOversellUnderConcurrency} 单测写死 \texttt{success == exactly 100}，少卖 / 超卖 CI 都红。本质是 Redis 单线程 + Lua 原子。代码：\texttt{inventory\_test.go:52--82}。
\end{frame}

\begin{frame}[allowframebreaks]{代码位置一览}
\footnotesize
\textbf{核心实现}
\begin{description}\setlength\itemsep{-1mm}
  \item[reserve / commit / release Lua]\srcnote{repository/cache/inventory.go:32--66}
  \item[InitStock + Snapshot]\srcnote{repository/cache/inventory.go:69--138}
  \item[SeedFromDB + 启动注入]\srcnote{service/inventory/syncer.go:14--48；initialize/inventory.go:11--17}
  \item[OrderCreate reserve + Saga 补偿]\srcnote{internal/order/service.go:101--155}
  \item[CancelUnpaidOrder release（幂等）]\srcnote{internal/order/cancel.go:14--67}
\end{description}
\textbf{关单口径（双口径自爆）}
\begin{description}\setlength\itemsep{-1mm}
  \item[RMQ TTL 30 分钟]\srcnote{repository/rabbitmq/order\_delay.go:22 OrderCancelDelay}
  \item[Cron 兜底 15 分钟]\srcnote{internal/order/task.go:24 GetTimeoutOrders(15, 100)}
\end{description}
\textbf{业务码（诚实标注）}
\begin{description}\setlength\itemsep{-1mm}
  \item[50001 = ErrorOss，非缺货码]\srcnote{pkg/e/code.go:49}
  \item[库存错误待映射业务码 90001]\srcnote{repository/cache/inventory.go:25 ErrStockInsufficient}
\end{description}
\textbf{测试 \& 压测}
\begin{description}\setlength\itemsep{-1mm}
  \item[NoOversell 500 抢 100 = 恰好 100]\srcnote{repository/cache/inventory\_test.go:52--82}
  \item[Lua vs DB 尾延迟 136ms vs 453ms]\srcnote{stressTest/REPORT.md \S 6}
\end{description}
\keypoint{配套业务 deck：09-order-lifecycle.tex / 11-consistency.tex / 12-merchant-ops.tex}
\end{frame}

\end{document}
